\documentclass[11pt,letterpaper]{exam}

\usepackage{amsmath}
\usepackage{amssymb}
\usepackage{graphicx}
\usepackage{geometry}
\usepackage{longtable}
\geometry{margin=1in}

\begin{document}

\begin{center}
\LARGE \textbf{Grading Rubric} \\
\large Midterm \\
\normalsize Duration: 75 minutes
\end{center}

\noindent \textbf{Instructions:} Answer all questions. Show your work for full credit. Use the formula sheet provided. You may use a calculator.

\vspace{0.25in}

\begin{questions}

\section*{Point Distribution}

\begin{longtable}{|p{0.5in}|p{4in}|r|}
\hline
\textbf{ID} & \textbf{Question} & \textbf{Points} \\
\hline
Q1 & What measure of central tendency is most affected by extreme... & 5 \\
\hline
Q2 & If a dataset has a standard deviation of 0, what can we conc... & 5 \\
\hline
Q3 & What is the probability of getting at least one head when fl... & 8 \\
\hline
Q4 & In hypothesis testing, a Type I error occurs when: & 8 \\
\hline
Q5 & Which of the following $p$-values provides the strongest evi... & 5 \\
\hline
Q6 & If two events $A$ and $B$ are independent, then $P(A \cap B)... & 8 \\
\hline
Q7 & Explain the difference between a population parameter and a ... & 12 \\
\hline
Q8 & A dataset has a mean of 50 and a standard deviation of 10. U... & 12 \\
\hline
Q9 & Explain what a $p$-value represents in the context of hypoth... & 15 \\
\hline
Q10 & A pharmaceutical company claims their new drug reduces blood... & 22 \\
\hline
\multicolumn{2}{|r|}{\textbf{Total Points:}} & \textbf{100} \\
\hline
\end{longtable}

\section*{Detailed Rubrics}

\subsection*{Q10: A pharmaceutical company claims their new drug red...}

Hypotheses stated correctly (6 pts), Type I error explained (5 pts), Type II error explained (5 pts), Analysis of error severity (6 pts)

\begin{itemize}
  \item \textbf{Hypotheses (6 pts)}: Correctly stated null hypothesis (μ = 0 or μ ≥ -10) and alternative hypothesis (μ < -10)
  \item \textbf{Type I Error (5 pts)}: Clearly explained as rejecting H₀ when true - drug approved when ineffective
  \item \textbf{Type II Error (5 pts)}: Clearly explained as failing to reject H₀ when false - drug not approved when effective
  \item \textbf{Error Analysis (6 pts)}: Thoughtful analysis of which error is more serious with valid reasoning about patient safety and resource allocation
\end{itemize}

\end{questions}

\end{document}
